\section{ماژول cur\_moe\_compressor (CUR-MoE Compressor)}
\label{sec:module_cur_moe_compressor}

\subsection{معرفی و انگیزه}
این مکانیزم بر مبنای مبانی نظری پیشرفته در ادبیات توجه و فشرده‌سازی \cite{ref_curmoecompressor_2024} توسعه یافته است.

افزایش مقیاس مدل‌های زبانی از طریق معماری‌های مبتنی بر مخلوطی از خبره‌ها (\lr{Mixture of Experts - MoE}) اگرچه امکان افزایش ظرفیت پارامتری بدون افزایش متناظر در هزینه محاسباتی هر توکن را فراهم می‌کند، اما چالش ذخیره‌سازی صدها میلیارد پارامتر در حافظه با پهنای باند بالا (\lr{HBM}) را به شدت تشدید می‌نماید. در لایه‌های خبره شبکه پیش‌خور (\lr{FFN})، ماتریس‌های پروجکشن مکرر با افزونگی اطلاعاتی مواجه هستند. روش‌های سنتی تجزیه مانند \lr{SVD} مستلزم تبدیل کامل پایه‌های متعامد بوده و تفسیرپذیری یا ساختار فیزیکی تنک را دگرگون می‌سازند؛ علاوه بر آن، هزینه محاسباتی تجزیه کامل مقادیر منفرد بر روی ده‌ها خبره مستقل به شدت بالا است.

ماژول فشرده‌ساز \lr{CUR-MoE Compressor} با تکیه بر تقریب متقاطع \lr{CUR Decomposition}، چارچوبی کارآمد برای فشرده‌سازی ساختاری لایه‌های خبره \lr{MoE} و نگاشت‌های چگال پیش‌خور ارائه می‌دهد. تقریب \lr{CUR} به جای پایه‌های انتزاعی تکین، ستون‌ها و سطرهای واقعی خود ماتریس وزن را بر مبنای امتیازات آماری نفوذ (\lr{Statistical Leverage Scores}) یا توزیع انرژی فروبنیوس (\lr{Frobenius Energy}) برمی‌گزیند. این راهبرد علاوه بر کاهش حافظه و تسریع استنتاج، ساختار ویژگی‌های خبره‌ها را حفظ کرده و تقریب شبه‌معکوس را با منظم‌سازی تیخونوف (\lr{Tikhonov-regularized Moore-Penrose Pseudoinverse}) به یک هسته فاکتورگیری پایدار تبدیل می‌نماید.

\subsection{مبانی نظری و فرمول‌بندی ریاضی}
فرض کنید ماتریس وزن یکی از لایه‌های خبره یا شبکه پیش‌خور به صورت $\mathbf{W} \in \mathbb{R}^{m \times n}$ تعریف شده باشد که در آن $m = d_{\text{out}}$ و $n = d_{\text{in}}$ است. تقریب \lr{CUR} به دنبال تجزیه این ماتریس به ضرب سه ماتریس با مرتبه تقلیل‌یافته $r \ll \min(m, n)$ است:
\begin{equation}
\mathbf{W} \approx \mathbf{C} \, \mathbf{U}_{\text{core}} \, \mathbf{R}
\end{equation}
که در آن $\mathbf{C} \in \mathbb{R}^{m \times c}$ زیرمجموعه‌ای شامل $c$ ستون از ماتریس $\mathbf{W}$، ماتریس $\mathbf{R} \in \mathbb{R}^{r_{\text{rows}} \times n}$ زیرمجموعه‌ای شامل $r_{\text{rows}}$ سطر از $\mathbf{W}$ (در حالت متقارن $c = r_{\text{rows}} = r$) و ماتریس $\mathbf{U}_{\text{core}} \in \mathbb{R}^{c \times r_{\text{rows}}}$ هسته میانی ارتباط‌دهنده است.

انتخاب سطرها و ستون‌ها بر اساس دو معیار اصلی انجام می‌پذیرد:
\begin{enumerate}
    \item \textbf{امتیازات نفوذ آماری (\lr{Leverage Scores}):} با محاسبه تجزیه مقادیر منفرد فشرده $\mathbf{W} = \mathbf{P} \mathbf{\Sigma} \mathbf{Q}^\top$، امتیاز نفوذ ستون $j$ و سطر $i$ به صورت زیر تعیین می‌شود:
    \begin{equation}
    p_j^{\text{col}} = \frac{1}{r} \sum_{k=1}^r (\mathbf{Q}_{j, k})^2, \quad p_i^{\text{row}} = \frac{1}{r} \sum_{k=1}^r (\mathbf{P}_{i, k})^2
    \end{equation}
    \item \textbf{انرژی فروبنیوس سطرها و ستون‌ها (\lr{Frobenius Energy}):} در روش بدون نیاز به \lr{SVD}، معیار نرم اقلیدسی بردارها به کار گرفته می‌شود:
    \begin{equation}
    e_j^{\text{col}} = \frac{\|\mathbf{W}_{:, j}\|_2^2}{\|\mathbf{W}\|_F^2}, \quad e_i^{\text{row}} = \frac{\|\mathbf{W}_{i, :}\|_2^2}{\|\mathbf{W}\|_F^2}
    \end{equation}
\end{enumerate}

پس از استخراج قطعی شاخص‌های برتر با سازوکار \lr{Top-k}، ماتریس اشتراک $\mathbf{W}_{\text{sub}} = \mathbf{W}_{I_{\text{rows}}, I_{\text{cols}}} \in \mathbb{R}^{r \times r}$ استخراج می‌شود. هسته میانی $\mathbf{U}_{\text{core}}$ از طریق شبه‌معکوس مور-پنروز منظم‌شده با تیخونوف با فاکتور میرایی $\lambda \ge 10^{-6}$ فرمول‌بندی می‌گردد:
\begin{equation}
\mathbf{U}_{\text{core}} = \mathbf{W}_{\text{sub}}^\dagger \approx (\mathbf{W}_{\text{sub}}^\top \mathbf{W}_{\text{sub}} + \lambda \mathbf{I})^{-1} \mathbf{W}_{\text{sub}}^\top
\end{equation}

تعداد کل پارامترهای لایه از $m \cdot n$ به $r(m + n + r)$ تقلیل می‌یابد. بنابراین قید کاهش پارامتری لایه فشرده به صورت زیر ارزیابی می‌شود:
\begin{equation}
r \cdot (m + n + r) < 0.85 \cdot (m \cdot n)
\end{equation}

\paragraph{تحلیل طیفی و نامساوی‌های پایداری خطا}
برای بررسی میزان دقت تقریب ماتریس وزن یا نگاشت توجه در ماژول \lr{cur\_moe\_compressor}، تفاضل ماتریس فاکتورگیری‌شده از نگاشت متراکم اصلی بر اساس نامساوی هافدینگ-فون‌نیومان ارزیابی می‌گردد. فرض کنید $\mathbf{W} \in \mathbb{R}^{d_{\text{in}} \times d_{\text{out}}}$ ماتریس وزن هدف باشد که توسط تقریب کم‌مرتبه یا پراکنده با رتبه موثر $r$ جایگزین شده است. کران بالای خطای بازسازی بر مبنای نرم فروبنیوس به صورت زیر تصریح می‌شود:
\begin{equation}
\label{eq:spectral_bound_cur_moe_compressor}
\| \mathbf{W} - \mathbf{W}_{r} \|_F = \sqrt{\sum_{i=r+1}^{\min(d_{\text{in}}, d_{\text{out}})} \sigma_i^2(\mathbf{W})} \le \epsilon \| \mathbf{W} \|_F
\end{equation}
که در آن $\sigma_i(\mathbf{W})$ بیانگر $i$-امین مقدار منفرد مرتب‌شده ماتریس به صورت نزولی است. نسبت حفظ انرژی طیفی $\eta_r = \frac{\sum_{i=1}^r \sigma_i^2}{\sum_{i=1}^{d} \sigma_i^2}$ تضمین می‌کند که با انتخاب آستانه $\tau \ge 0.90$، بیش از نود درصد از پراکندگی واریانس داده‌ها حفظ شده و افت سیگنال در انتشار گرادیان رو به عقب به حداقل ممکن کاهش می‌یابد.

\paragraph{تحلیل مجانبی و پیچیدگی اشتراک حافظه سخت‌افزاری}
از منظر پیچیدگی محاسباتی و انتقال داده‌ها در سلسومراتب حافظه (حافظه سراسری \lr{HBM} به حافظه محلی \lr{SRAM})، عملیات ضرب ماتریسی در این مکانیزم نیازمند تخصیص بافر موقت با ابعاد فاکتورگیری‌شده است. برای پردازش دنباله‌ای به طول $N$ با بعد مخفی $D$، بار محاسباتی هسته از فرم درجه دو به رابطه زیر تقلیل می‌یابد:
\begin{equation}
\label{eq:complexity_cur_moe_compressor}
\mathcal{C}_{\text{cur\_moe\_compressor}} = \mathcal{O}\left( B \cdot N \cdot D \cdot r + B \cdot H \cdot N \cdot d_k \right) \ll \mathcal{O}\left( B \cdot H \cdot N^2 \cdot d_k \right)
\end{equation}
این بهبود مرتبه پیچیدگی، با کاهش چشمگیر حجم تبادل تانسورها در گذرگاه‌های ارتباطی کارت گرافیک، مانع از اشباع پهنای باند حافظه (پدیده \lr{Memory-Bound}) در فرآیند استنتاج طولانی‌مدت می‌گردد.

\subsection{معماری و پیاده‌سازی}
کلاس \lr{CURMoECompressor} در ماژول \lr{Codes/framework/compression\_engine/cur\_moe\_compressor.py} پیاده‌سازی شده و تحت شناسه‌های \lr{cur\_moe}، \lr{cur\_moe\_class} و \lr{cur\_moe\_compressor} در رجیستری فشرده‌سازی ثبت گردیده است. مؤلفه‌های کلیدی معماری عبارتند از:
\begin{itemize}
    \item \textbf{هدف‌گیری اختصاصی خبره‌ها (\lr{MoE \& FFN Targeting}):} این ماژول دارای پارامتر \lr{target\_attention\_only} است که در صورت غیرفعال بودن، به‌طور خاص روی لایه‌های فیدفوروارد و خبره‌های پرحجم متمرکز می‌شود تا حافظه لایه‌های \lr{Intermediate / MLP} را آزاد کند.
    \item \textbf{لایه‌های خطی فاکتورگیری‌شده \lr{CURLinear}:} لایه‌های چگال استاندارد و لایه‌های \lr{Conv1D} (در مدل‌های دیکودر) با ماژول دو مرحله‌ای جایگزین می‌شوند که ضرب پیشرو را به صورت $\mathbf{y} = ((\mathbf{x} \mathbf{R}^\top) \mathbf{U}_{\text{core}}^\top) \mathbf{C}^\top + \mathbf{b}$ اجرا می‌کند.
    \item \textbf{پلکان تطبیقی رتبه و خطای فروبنیوس (\lr{Adaptive Rank Escalation}):} پس از تقریب لایه، خطای نسبی بازسازی فروبنیوس محاسبه می‌شود. در صورتی که خطا از آستانه مجاز (\lr{max\_error=0.90}) فراتر رود، الگوریتم به صورت خودکار تا دو بار (\lr{max\_retries=2}) اقدام به ارتقای رتبه نموده و در صورت تداوم خطا، با سازوکار بازگشت ایمن (\lr{Safe Fallback})، وزن متراکم دست‌نخورده حفظ می‌گردد.
\end{itemize}

\subsection{نتایج تجربی}
عملکرد فشرده‌ساز \lr{cur\_moe\_compressor} بر اساس اطلاعات بنچ‌مارک رسمی در \lr{thesis\_final/master\_verified\_results.json} روی مدل‌های پایه به شرح جدول \ref{tab:benchmark_cur_moe_compressor} استخراج شده است:

\begin{table}[H]
\centering
\small
\caption{نتایج بنچ‌مارک تجربی ماژول فشرده‌ساز \lr{cur\_moe\_compressor} روی مدل‌های زبانی پایه}
\label{tab:benchmark_cur_moe_compressor}
\begin{tabular}{lccccc}
\toprule
\textbf{مدل پایه \lr{(Base Model)}} & \textbf{دقت \lr{(Accuracy)}} & \textbf{شاخص \lr{F1}} & \textbf{حافظه \lr{(Peak VRAM)}} & \textbf{توان \lr{(Throughput)}} & \textbf{وضعیت} \\
\midrule
\lr{DistilBERT} & \lr{50.92\%} & \lr{0.6748} & \lr{3611.10 MB} & \lr{117,166.70 tok/s} & \lr{PASS} \\
\lr{RoBERTa} & \lr{50.34\%} & \lr{---} & \lr{---} & \lr{---} & \lr{SKIP\_ROBERTA\_ARCHITECTURAL} \\
\lr{GPT-2} & \lr{74.66\%} & \lr{---} & \lr{---} & \lr{---} & \lr{PASS} \\
\lr{Qwen} & \lr{52.71\%} & \lr{---} & \lr{3205.92 MB} & \lr{---} & \lr{PASS} \\
\bottomrule
\end{tabular}
\end{table}

بر اساس نتایج جدول \ref{tab:benchmark_cur_moe_compressor}، ماژول \lr{cur\_moe\_compressor} توانسته در مدل دیکودر \lr{GPT-2} به دقت \lr{74.66\%} دست یابد. در مدل \lr{DistilBERT} با ثبت توان پردازشی بالای \lr{117,166.70 tok/s} وضعیت \lr{PASS} را احراز کرده است، هرچند دقت دسته‌بندی \lr{SST-2} در سطح شانس تصادفی (\lr{50.92\%}) متوقف شده است. در مدل \lr{RoBERTa}، به دلیل تداخل تزریق لایه‌های متناظر خبره با ساختار ماسک پیش‌آموزش و ناسازگاری توابع فعال‌سازی، وضعیت آزمایش به صورت مستند به عنوان \lr{SKIP\_ROBERTA\_ARCHITECTURAL} ثبت شده است. در مدل زبانی بزرگ \lr{Qwen} نیز دقت \lr{52.71\%} با حافظه فعال \lr{3205.92 MB} حاصل گردیده است. همچنین در ارزیابی خط مبنای مجزای \lr{cur\_moe}، دقت \lr{67.43\%} در \lr{DistilBERT} با برچسب \lr{PASS} و وضعیت‌های \lr{PASS\_LOW} در مدل‌های زبانی علّی ثبت شده است.

\subsection{محدودیت‌ها و تحلیل شکست}
تحلیل ساختاری رفتار ماژول \lr{cur\_moe\_compressor} در آزمایش‌های تجربی، نقاط ضعف و سناریوهای شکست زیر را آشکار می‌سازد:
\begin{itemize}
    \item \textbf{رانش توزیع در مدل‌های بدون ساختار صریح خبره (\lr{Dense FFN Drift}):} استفاده از تقریب \lr{CUR} برای لایه‌های چگال پیش‌خور متوالی در مدل‌های متراکم انکودر (\lr{DistilBERT} و \lr{RoBERTa})، به دلیل فقدان روتر فعال‌ساز پویا، منجر به تضعیف خطوط غیرخطی \lr{GELU} و افت دقت تا سطح شانس (\lr{$\sim$50.9\%}) شده است.
    \item \textbf{بدحالی ماتریس میانی (\lr{Ill-Conditioned Core Matrix}):} در صورتی که سطرهای انتخاب‌شده با رتبه تنک نمونه‌برداری شوند، ماتریس تقاطع $\mathbf{W}_{\text{sub}}$ با مقادیر تکین نزدیک به صفر مواجه شده و شبه‌معکوس تیخونوف خطای بازسازی فروبنیوس را افزایش می‌دهد.
    \item \textbf{حساسیت بالا به عدم تطابق معماری در \lr{RoBERTa}:} ماژول‌های پیش‌خور انکودر در روبرتا نیازمند تنظیم ظریف‌تر ضرایب مقیاس هستند و فاکتورگیری ایستا بدون بازیابی گرادیانی کافی منجر به انحراف فضای ویژگی‌ها می‌گردد.
\end{itemize}
